% !TeX root = 20261022_exp87.tex
% !LW recipe=lablog XeLaTeX
\input{_support/style.tex}
\begin{document}
\reporttitle{2026-10-22}{EXP87 · 극소수 공격 사례의 TabPFN·XGB 분류 비교 — EXP86 학습곡선 확장}
\subsection{목적과 실행 상태}
EXP87은 EXP86의 공격-only 학습곡선에서 공격당 학습 사례 수를 20개보다 크게 늘린 후속 실험이다. 같은 모델·같은 공격 test·같은 seed를 유지하여 표본 수에 따른 Macro-F1과 클래스별 품질 변화를 확인한다. XGB n=\allowbreak{}1의 분할 불가 조건을 제외하여 2개부터 시작하며, 정상 포함 평가는 사용자 요청에 따라 제외한다.

현재 상태는 complete, 완료 조건은 200/\allowbreak{}200개다. 전체 x축에서 seed 43–47의 5회 평균·표본 표준편차를 비교한다. 5개 seed가 모두 끝나지 않은 조건은 주 곡선에서 제외하고 그림에 대기 중인 n을 명시한다.

\subsection{공통 설정과 평가 의미}
XGB는 기존 300 trees/\allowbreak{}depth8/\allowbreak{}lr0.05/\allowbreak{}subsample0.8/\allowbreak{}colsample0.8/\allowbreak{}min\_\allowbreak{}child\_\allowbreak{}weight1을 유지한다. PFN은 기존 v3 checkpoint, 4 estimators, fit\_\allowbreak{}with\_\allowbreak{}cache, CPU cache이며 가중치 FT와 context subsampling을 하지 않는다. 초기 1000개 context를 사용하지 않으며 각 공격당 정확히 n개씩 제공한다. 같은 seed·n에서 두 모델의 입력 행 ID가 같다.

EXP86에서 사용한 seed별 각 공격의 첫 20행을 그대로 유지하고, 같은 클래스의 남은 정제 train에서 중복 없이 무작위 추출하여 뒤에 붙인다. 따라서 작은 n의 표본은 큰 n에 포함되며, 정상 없이 균형 context를 구성한다. CIC 공격 6종, ToN 공격 9종이므로 전체 context는 각각 6n, 9n행이다.

평가는 EXP86 Phase A와 동일한 고정 공격 test의 Macro-F1이며 각 공격의 F1을 같은 가중치로 평균한다. 각 seed에서 계산한 점수를 평균하며 seed 간 혼동행렬을 먼저 합치지 않는다. 정상 test는 추론에서도 제외하고 attack-only query 순서와 batch8192를 모든 n에서 동일하게 사용한다. 특징 전처리와 모델 설정은 EXP86을 유지한다.

\subsection{실제 제공·평가 표본 수}
\begingroup
\footnotesize
\begin{longtable}{>{\raggedright\arraybackslash}p{\dimexpr 0.14861\linewidth-2\tabcolsep\relax}>{\raggedleft\arraybackslash}p{\dimexpr 0.14861\linewidth-2\tabcolsep\relax}>{\raggedright\arraybackslash}p{\dimexpr 0.35694\linewidth-2\tabcolsep\relax}>{\raggedright\arraybackslash}p{\dimexpr 0.19722\linewidth-2\tabcolsep\relax}>{\raggedleft\arraybackslash}p{\dimexpr 0.14861\linewidth-2\tabcolsep\relax}}
\toprule
데이터 & 공격 수 & 공격당 n & 전체 학습/\allowbreak{}context & 공격 test \\ \midrule
\endfirsthead
데이터 & 공격 수 & 공격당 n & 전체 학습/\allowbreak{}context & 공격 test \\ \midrule
\endhead
{}CIC2018 & 6 & 2/\allowbreak{}5/\allowbreak{}10/\allowbreak{}20/\allowbreak{}40/\allowbreak{}80/\allowbreak{}160/\allowbreak{}320/\allowbreak{}450 & 6n; 최대 2,700행 & 390,270 \\
{}ToN-IoT & 9 & 2/\allowbreak{}5/\allowbreak{}10/\allowbreak{}20/\allowbreak{}40/\allowbreak{}80/\allowbreak{}160/\allowbreak{}320/\allowbreak{}640/\allowbreak{}1280/\allowbreak{}2151 & 9n; 최대 19,359행 & 797,421 \\
\bottomrule
\end{longtable}
\endgroup
가로축 끝점은 성능 포화가 아니라 균등 표본 선택의 데이터 상한이다. CIC2018의 web\_\allowbreak{}attacks 학습 사례는 450개, ToN-IoT의 ransomware는 2,151개여서 모든 공격에 같은 n을 주는 현재 조건에서는 더 늘릴 수 없다.

\subsection{Attack-only Macro-F1 (\%)}
EXP86과 같이 XGB는 검정색, TabPFN은 파란색이며 점은 5-seed 평균, 오차 막대는 표본 표준편차다. 각 점에 평균을 소수 둘째 자리까지 표시하고 아래 패널에는 TabPFN−XGB의 차이(\%p)를 표시한다. 2개부터 수천 개까지 한 곡선에서 읽을 수 있도록 x축은 로그 간격이며 실제 n을 모두 표시한다. 모든 조건의 평가 집합과 평균에 포함한 seed는 동일하다.

\reportfigure{figures/exp87_quality_cic2018.pdf}{CIC2018: 5 seeds 평균±표본 표준편차; 점 옆 수치는 평균, 아래 패널은 모델 간 차이}


\reportfigure{figures/exp87_quality_toniot.pdf}{ToN-IoT: 5 seeds 평균±표본 표준편차; 점 옆 수치는 평균, 아래 패널은 모델 간 차이}


\subsection{수치 요약}
두 모델 모두 동일 5개 seed가 완료된 조건의 평균±표본 표준편차다. 모든 원시 결과와 부분 완료 조건은 values.csv와 summary.csv에 보존한다. 표준편차는 선택 표본·모델 seed 변동이며 신뢰구간이나 통계적 동등성을 뜻하지 않는다.

\begingroup
\footnotesize
\begin{longtable}{>{\raggedright\arraybackslash}p{\dimexpr 0.16604\linewidth-2\tabcolsep\relax}>{\raggedleft\arraybackslash}p{\dimexpr 0.16604\linewidth-2\tabcolsep\relax}>{\raggedright\arraybackslash}p{\dimexpr 0.21321\linewidth-2\tabcolsep\relax}>{\raggedright\arraybackslash}p{\dimexpr 0.23208\linewidth-2\tabcolsep\relax}>{\raggedright\arraybackslash}p{\dimexpr 0.22264\linewidth-2\tabcolsep\relax}}
\toprule
데이터 & n & XGB F1 (\%) & TabPFN F1 (\%) & PFN−XGB (\%p) \\ \midrule
\endfirsthead
데이터 & n & XGB F1 (\%) & TabPFN F1 (\%) & PFN−XGB (\%p) \\ \midrule
\endhead
{}CIC2018 & 2 & 63.24 ± 6.32 & 66.34 ± 10.44 & +\allowbreak{}3.10 \\
{}CIC2018 & 5 & 72.24 ± 4.66 & 74.78 ± 4.29 & +\allowbreak{}2.55 \\
{}CIC2018 & 10 & 80.37 ± 1.91 & 80.58 ± 0.73 & +\allowbreak{}0.21 \\
{}CIC2018 & 20 & 80.43 ± 0.97 & 80.59 ± 1.64 & +\allowbreak{}0.16 \\
{}CIC2018 & 40 & 81.25 ± 1.08 & 81.78 ± 1.27 & +\allowbreak{}0.54 \\
{}CIC2018 & 80 & 82.22 ± 1.19 & 83.96 ± 0.49 & +\allowbreak{}1.75 \\
{}CIC2018 & 160 & 85.30 ± 3.28 & 84.79 ± 0.55 & -0.50 \\
{}CIC2018 & 320 & 86.39 ± 4.19 & 84.78 ± 0.58 & -1.60 \\
{}CIC2018 & 450 & 86.25 ± 3.81 & 85.42 ± 0.79 & -0.83 \\
{}ToN-IoT & 2 & 45.03 ± 3.27 & 45.77 ± 6.04 & +\allowbreak{}0.74 \\
{}ToN-IoT & 5 & 52.29 ± 5.62 & 58.75 ± 3.67 & +\allowbreak{}6.46 \\
{}ToN-IoT & 10 & 57.88 ± 5.08 & 60.46 ± 3.78 & +\allowbreak{}2.58 \\
{}ToN-IoT & 20 & 64.74 ± 2.18 & 69.13 ± 4.49 & +\allowbreak{}4.39 \\
{}ToN-IoT & 40 & 72.70 ± 4.20 & 75.46 ± 3.37 & +\allowbreak{}2.75 \\
{}ToN-IoT & 80 & 78.26 ± 0.75 & 81.28 ± 4.59 & +\allowbreak{}3.02 \\
{}ToN-IoT & 160 & 82.76 ± 0.94 & 84.95 ± 2.59 & +\allowbreak{}2.19 \\
{}ToN-IoT & 320 & 86.02 ± 1.63 & 86.96 ± 2.02 & +\allowbreak{}0.93 \\
{}ToN-IoT & 640 & 89.82 ± 2.27 & 88.71 ± 1.46 & -1.10 \\
{}ToN-IoT & 1,280 & 91.83 ± 0.35 & 88.17 ± 2.12 & -3.66 \\
{}ToN-IoT & 2,151 & 92.75 ± 0.45 & 91.12 ± 2.41 & -1.63 \\
\bottomrule
\end{longtable}
\endgroup
\subsection{클래스별 결과}
EXP86의 클래스별 F1 그림을 같은 n 범위까지 확장한다. 각 셀은 5-seed 평균이며 클래스는 고정 공격 test 표본 수 내림차순이다. 반복이 남은 열은 공란으로 표시한다.

\reportfigure{figures/exp87_per_class_cic2018.pdf}{CIC2018: 클래스별 F1의 5-seed 평균; 고정 공격 test 표본 수 내림차순}


\reportfigure{figures/exp87_per_class_toniot.pdf}{ToN-IoT: 클래스별 F1의 5-seed 평균; 고정 공격 test 표본 수 내림차순}


\subsection{비용}
EXP86과 같이 초기 준비와 추론을 분리한다. 준비는 새 모델 생성과 XGB fit 또는 PFN context/\allowbreak{}cache 준비의 합이며 순차 갱신 비용이 아니다. 이번 추론시간은 공격 test만 처리한 시간이므로 EXP86의 정상 포함 전체 test 추론시간과 직접 비교하지 않는다.

\reportfigure{figures/exp87_cost.pdf}{초기 준비와 공격 test 추론시간; 동일 5 seeds 평균±표본 표준편차}


\subsection{근거에 따른 해석과 추가 확대 가능성}
현재 결과는 수렴을 확정하지 않는다. 이전 보고서는 수렴 절차를 앞세우고 단일 seed 곡선과 마지막 세 조건의 평균 곡선을 분리하여 EXP86에서 이어지는 학습곡선을 읽기 어렵게 했다. 이번 개정에서는 전체 범위를 동일 5-seed 평균으로 연결하고 수렴 규칙과 실행 이력을 뒤로 옮겼다.

CIC2018에서 평균 차이가 1\%p 이내인 측정 조건은 n=\allowbreak{}[10, 20, 40, 160, 450]다. 이는 점수 차이를 읽기 위한 기술적 기준이며 통계적 동등성이나 그 뒤의 모든 n에서 비슷함을 의미하지 않는다.

CIC2018 XGB의 마지막 두 완료 조건은 n=\allowbreak{}320 \ensuremath{\rightarrow} 450, 평균 F1 86.39 \ensuremath{\rightarrow} 86.25\% (변화 -0.14\%p)다.

CIC2018 TabPFN의 마지막 두 완료 조건은 n=\allowbreak{}320 \ensuremath{\rightarrow} 450, 평균 F1 84.78 \ensuremath{\rightarrow} 85.42\% (변화 +\allowbreak{}0.64\%p)다.

ToN-IoT에서 평균 차이가 1\%p 이내인 측정 조건은 n=\allowbreak{}[2, 320]다. 이는 점수 차이를 읽기 위한 기술적 기준이며 통계적 동등성이나 그 뒤의 모든 n에서 비슷함을 의미하지 않는다.

ToN-IoT XGB의 마지막 두 완료 조건은 n=\allowbreak{}1,280 \ensuremath{\rightarrow} 2,151, 평균 F1 91.83 \ensuremath{\rightarrow} 92.75\% (변화 +\allowbreak{}0.93\%p)다.

ToN-IoT TabPFN의 마지막 두 완료 조건은 n=\allowbreak{}1,280 \ensuremath{\rightarrow} 2,151, 평균 F1 88.17 \ensuremath{\rightarrow} 91.12\% (변화 +\allowbreak{}2.96\%p)다.

특히 ToN-IoT에서 마지막 구간의 개선이 계속되어 더 큰 학습 집합을 시험할 근거가 있다. CIC2018에서는 전체 평균 변화가 작더라도 클래스별 품질과 seed 변동을 함께 확인해야 한다. 서로 다른 공격 정의·클래스 수·특징 분포를 가진 두 데이터셋의 고유 난이도를 점수만으로 단정하지 않는다.

균등한 공격당 n을 더 늘리려면 가장 작은 공격 클래스의 새로운 독립 학습 사례가 필요하다. 현재 데이터만으로는 희귀 공격을 모두 유지하고 나머지 공격의 사례를 늘리는 대안을 구성할 수 있다. 후자는 클래스 비율도 달라지므로 기존 균형 곡선과 구분하여 보고해야 하며, 같은 행의 복제를 새로운 사례로 세지 않는다.

\subsection{마지막 측정점의 precision·recall·F1}
각 데이터에서 현재 가장 큰 완료 n의 수치이며, 가능한 seed만 평균한다. 낮은 F1이 누락 때문인지 다른 공격의 false positive 유입 때문인지 precision·recall을 함께 본다. 행은 고정 공격 test 표본 수 내림차순이다.

\begingroup
\footnotesize
\begin{longtable}{>{\raggedright\arraybackslash}p{\dimexpr 0.16298\linewidth-2\tabcolsep\relax}>{\raggedright\arraybackslash}p{\dimexpr 0.15594\linewidth-2\tabcolsep\relax}>{\raggedright\arraybackslash}p{\dimexpr 0.12072\linewidth-2\tabcolsep\relax}>{\raggedleft\arraybackslash}p{\dimexpr 0.12072\linewidth-2\tabcolsep\relax}>{\raggedleft\arraybackslash}p{\dimexpr 0.17002\linewidth-2\tabcolsep\relax}>{\raggedleft\arraybackslash}p{\dimexpr 0.14889\linewidth-2\tabcolsep\relax}>{\raggedleft\arraybackslash}p{\dimexpr 0.12072\linewidth-2\tabcolsep\relax}}
\toprule
데이터 /\allowbreak{} n & 클래스 & 모델 & seeds & Precision (\%) & Recall (\%) & F1 (\%) \\ \midrule
\endfirsthead
데이터 /\allowbreak{} n & 클래스 & 모델 & seeds & Precision (\%) & Recall (\%) & F1 (\%) \\ \midrule
\endhead
{}cic2018 /\allowbreak{} 450 & ddos & xgb & 5 & 100.00 & 99.23 & 99.61 \\
{}cic2018 /\allowbreak{} 450 & ddos & pfn & 5 & 100.00 & 99.66 & 99.83 \\
{}cic2018 /\allowbreak{} 450 & bot & xgb & 5 & 99.91 & 98.61 & 99.25 \\
{}cic2018 /\allowbreak{} 450 & bot & pfn & 5 & 99.99 & 99.13 & 99.56 \\
{}cic2018 /\allowbreak{} 450 & dos & xgb & 5 & 99.99 & 100.00 & 100.00 \\
{}cic2018 /\allowbreak{} 450 & dos & pfn & 5 & 100.00 & 98.37 & 99.18 \\
{}cic2018 /\allowbreak{} 450 & brute\_\allowbreak{}force & xgb & 5 & 100.00 & 99.95 & 99.97 \\
{}cic2018 /\allowbreak{} 450 & brute\_\allowbreak{}force & pfn & 5 & 100.00 & 99.95 & 99.97 \\
{}cic2018 /\allowbreak{} 450 & infiltration & xgb & 5 & 86.10 & 97.11 & 91.17 \\
{}cic2018 /\allowbreak{} 450 & infiltration & pfn & 5 & 88.25 & 97.56 & 92.54 \\
{}cic2018 /\allowbreak{} 450 & web\_\allowbreak{}attacks & xgb & 5 & 17.48 & 97.64 & 27.48 \\
{}cic2018 /\allowbreak{} 450 & web\_\allowbreak{}attacks & pfn & 5 & 12.03 & 99.53 & 21.44 \\
{}toniot /\allowbreak{} 2151 & xss & xgb & 5 & 99.72 & 97.79 & 98.74 \\
{}toniot /\allowbreak{} 2151 & xss & pfn & 5 & 99.78 & 97.40 & 98.57 \\
{}toniot /\allowbreak{} 2151 & password & xgb & 5 & 99.76 & 99.70 & 99.73 \\
{}toniot /\allowbreak{} 2151 & password & pfn & 5 & 99.81 & 99.62 & 99.72 \\
{}toniot /\allowbreak{} 2151 & ddos & xgb & 5 & 99.41 & 97.75 & 98.58 \\
{}toniot /\allowbreak{} 2151 & ddos & pfn & 5 & 99.30 & 98.21 & 98.75 \\
{}toniot /\allowbreak{} 2151 & injection & xgb & 5 & 92.84 & 97.02 & 94.88 \\
{}toniot /\allowbreak{} 2151 & injection & pfn & 5 & 92.59 & 97.20 & 94.84 \\
{}toniot /\allowbreak{} 2151 & backdoor & xgb & 5 & 99.98 & 99.92 & 99.95 \\
{}toniot /\allowbreak{} 2151 & backdoor & pfn & 5 & 99.98 & 99.95 & 99.96 \\
{}toniot /\allowbreak{} 2151 & dos & xgb & 5 & 86.79 & 99.10 & 92.53 \\
{}toniot /\allowbreak{} 2151 & dos & pfn & 5 & 92.93 & 98.47 & 95.47 \\
{}toniot /\allowbreak{} 2151 & scanning & xgb & 5 & 91.07 & 98.18 & 94.47 \\
{}toniot /\allowbreak{} 2151 & scanning & pfn & 5 & 87.13 & 98.50 & 92.28 \\
{}toniot /\allowbreak{} 2151 & mitm & xgb & 5 & 40.00 & 98.52 & 56.83 \\
{}toniot /\allowbreak{} 2151 & mitm & pfn & 5 & 28.97 & 99.22 & 43.26 \\
{}toniot /\allowbreak{} 2151 & ransomware & xgb & 5 & 98.33 & 99.83 & 99.08 \\
{}toniot /\allowbreak{} 2151 & ransomware & pfn & 5 & 94.96 & 99.72 & 97.28 \\
\bottomrule
\end{longtable}
\endgroup
\subsection{재현 점검과 실행 이력}
처음에는 seed43으로 전체 범위를 탐색하고 마지막 세 n만 seed44–47로 확인했다. 사용자 피드백에 따라 중간 구간까지 나머지 seed를 채워 모든 n에서 같은 반복 수로 비교한다. 완료한 기존 88조건은 원본 결과를 보존하여 재사용하고 빠진 112조건만 추가한다.

현재 상태는 complete /\allowbreak{} complete, 완료 모델 조건은 200개다. 실행 시작은 2026-10-22T16:37:40.334640+\allowbreak{}09:00, 기록 시각은 2026-10-22T17:05:21.168992+\allowbreak{}09:00다. 실행은 v7/\allowbreak{}EXP87에 등록했으며 경로는 \texttt{tabpfn/\allowbreak{}\allowbreak{}r\allowbreak{}esults/\allowbreak{}\allowbreak{}v\allowbreak{}7/\allowbreak{}\allowbreak{}exp87/\allowbreak{}\allowbreak{}20261022\allowbreak{}\_\allowbreak{}\allowbreak{}163739\_\allowbreak{}\allowbreak{}full\_\allowbreak{}\allowbreak{}lea\allowbreak{}rning\_\allowbreak{}\allowbreak{}cu\allowbreak{}rve\_\allowbreak{}\allowbreak{}five\allowbreak{}\_\allowbreak{}\allowbreak{}seeds\_\allowbreak{}\allowbreak{}c\allowbreak{}59213f5}다.

cic2018: seed43 측정 n=\allowbreak{}[2, 5, 10, 20, 40, 80, 160, 320, 450], 단일 seed 후보 확인=\allowbreak{}False, 다중 seed 상태=\allowbreak{}complete, 라벨 상한 도달=\allowbreak{}True.

최근 확인 구간 [160, 320, 450], 사용 seed [43, 44, 45, 46, 47]: XGB 점수 범위 1.09\%p, PFN 점수 범위 0.64\%p. 해당 범위는 최대−최소 차이다.

가용 라벨 상한까지 평가했지만 5-seed 평균의 포화 기준은 만족하지 않았다. 더 큰 균형 n은 현재 train pool에서 구성할 수 없다.

toniot: seed43 측정 n=\allowbreak{}[2, 5, 10, 20, 40, 80, 160, 320, 640, 1280, 2151], 단일 seed 후보 확인=\allowbreak{}False, 다중 seed 상태=\allowbreak{}complete, 라벨 상한 도달=\allowbreak{}True.

최근 확인 구간 [640, 1280, 2151], 사용 seed [43, 44, 45, 46, 47]: XGB 점수 범위 2.94\%p, PFN 점수 범위 2.96\%p. 해당 범위는 최대−최소 차이다.

가용 라벨 상한까지 평가했지만 5-seed 평균의 포화 기준은 만족하지 않았다. 더 큰 균형 n은 현재 train pool에서 구성할 수 없다.

최초 탐색은 최근 세 n의 점수 범위가 두 모델 모두 1\%p 이내인지 확인하는 방식이었다. 실제 종료 원인은 데이터 상한이며 수렴 확정이 아니다. 모델 간 차이가 1\%p 이내인 것과 각 모델이 더 이상 개선되지 않는 것은 다르다.

\subsection{EXP86과 query 배치 변경 점검}
추가로 n=\allowbreak{}2/\allowbreak{}5/\allowbreak{}10/\allowbreak{}20의 두 모델·5 seeds 전체 80조건을 EXP86과 대조했다. 모든 학습 행 ID가 동일하며, 정상 query를 제외한 배치 변경 후 Macro-F1 절대 차이의 최댓값은 0.00686\%p다. 해당 구간의 연속성을 exp86\_\allowbreak{}continuity\_\allowbreak{}check.json에 보존한다.

동일한 n=\allowbreak{}20 학습 행·공격 test에서 정상 query를 제거한 실행과 EXP86의 예측을 비교한다. 모든 EXP87 n은 공격만 입력하는 방식으로 실행했다. 아래 비교는 query 구성·배치 경계 변경에 대한 진단이며 기존 전체 test 추론시간과 새 공격-only 추론시간은 평가 행 수가 다르다.

cic2018 seed43 xgb: 예측 불일치 0행 (0.0000\%), Macro-F1 변화 0.0000\%p.

cic2018 seed43 pfn: 예측 불일치 0행 (0.0000\%), Macro-F1 변화 0.0000\%p.

toniot seed43 xgb: 예측 불일치 0행 (0.0000\%), Macro-F1 변화 0.0000\%p.

toniot seed43 pfn: 예측 불일치 2행 (0.0003\%), Macro-F1 변화 -0.0001\%p.

cic2018 seed44 xgb: 예측 불일치 0행 (0.0000\%), Macro-F1 변화 0.0000\%p.

cic2018 seed44 pfn: 예측 불일치 0행 (0.0000\%), Macro-F1 변화 0.0000\%p.

cic2018 seed45 xgb: 예측 불일치 0행 (0.0000\%), Macro-F1 변화 0.0000\%p.

cic2018 seed45 pfn: 예측 불일치 0행 (0.0000\%), Macro-F1 변화 0.0000\%p.

cic2018 seed46 xgb: 예측 불일치 0행 (0.0000\%), Macro-F1 변화 0.0000\%p.

cic2018 seed46 pfn: 예측 불일치 0행 (0.0000\%), Macro-F1 변화 0.0000\%p.

cic2018 seed47 xgb: 예측 불일치 0행 (0.0000\%), Macro-F1 변화 0.0000\%p.

cic2018 seed47 pfn: 예측 불일치 0행 (0.0000\%), Macro-F1 변화 0.0000\%p.

toniot seed44 xgb: 예측 불일치 0행 (0.0000\%), Macro-F1 변화 0.0000\%p.

toniot seed44 pfn: 예측 불일치 1행 (0.0001\%), Macro-F1 변화 -0.0003\%p.

toniot seed45 xgb: 예측 불일치 0행 (0.0000\%), Macro-F1 변화 0.0000\%p.

toniot seed45 pfn: 예측 불일치 1행 (0.0001\%), Macro-F1 변화 -0.0001\%p.

toniot seed46 xgb: 예측 불일치 0행 (0.0000\%), Macro-F1 변화 0.0000\%p.

toniot seed46 pfn: 예측 불일치 2행 (0.0003\%), Macro-F1 변화 0.0001\%p.

toniot seed47 xgb: 예측 불일치 0행 (0.0000\%), Macro-F1 변화 0.0000\%p.

toniot seed47 pfn: 예측 불일치 1행 (0.0001\%), Macro-F1 변화 0.0000\%p.

\subsection{보조 평가의 범위}
EXP86의 feature-disjoint attack은 해당 학습/\allowbreak{}context에 포함한 특징 벡터와 완전히 동일한 공격 test 행을 제외한 보조 집계다. 서로 다른 원본 행 ID라도 숫자 특징 전체가 같을 수 있다. 이는 다른 특징 열을 사용한다는 뜻이나 세션·근접 중복까지 제거했다는 뜻이 아니다. 제외 행이 n·seed마다 달라지므로 EXP87의 주 결과와 포화 판단에는 사용하지 않고 고정 attack-only test를 유지한다. 동일 특징 중첩 수는 결과 JSON/\allowbreak{}CSV에 진단 자료로 남긴다.

이미 검토한 고정 test의 탐색적 학습곡선이며 추가 seed는 표본 선택·모델 변동을 확인한다. 이 실험만으로 새로운 네트워크 환경에서의 일반화나 전체 IDS 운영 성능을 입증하지 않는다.

\subsection{추가 확대: 희귀 공격은 전부 유지하고 나머지 공격의 사례 수 증가}
사용자가 추가 확대안을 승인하여 seed43으로 실행한다. 현재 상태는 complete /\allowbreak{} complete, 추가 조건은 10/\allowbreak{}10개 완료다. 위의 5-seed 균형 곡선과 달리 아래는 seed43의 별도 조건이다.

각 클래스의 사용 수는 min(목표 n, 해당 클래스의 가용 train 수)다. 기존 균형 조건에서 사용한 모든 사례를 유지하면서 나머지 서로 다른 행을 추가하고 같은 고정 공격 test를 평가한다. 정상·복제 행·가중치 fine-tuning은 추가하지 않는다. 표본 수와 클래스 비율이 함께 변하므로 개선이 나타나더라도 순수한 표본 수 효과로 단정하지 않는다.

CIC2018의 실제 클래스별 제공 수다. 첫 열은 기존 균형 조건의 seed43 연결점(n=\allowbreak{}450)이며, 행은 고정 공격 test 표본 수 내림차순이다.

\begingroup
\footnotesize
\begin{longtable}{>{\raggedright\arraybackslash}p{\dimexpr 0.25854\linewidth-2\tabcolsep\relax}>{\raggedleft\arraybackslash}p{\dimexpr 0.18537\linewidth-2\tabcolsep\relax}>{\raggedleft\arraybackslash}p{\dimexpr 0.18537\linewidth-2\tabcolsep\relax}>{\raggedleft\arraybackslash}p{\dimexpr 0.18537\linewidth-2\tabcolsep\relax}>{\raggedleft\arraybackslash}p{\dimexpr 0.18537\linewidth-2\tabcolsep\relax}}
\toprule
공격 클래스 & 기존 450 & 800 & 1600 & 3200 \\ \midrule
\endfirsthead
공격 클래스 & 기존 450 & 800 & 1600 & 3200 \\ \midrule
\endhead
{}ddos & 450 & 800 & 1,600 & 3,200 \\
{}bot & 450 & 800 & 1,600 & 3,200 \\
{}dos & 450 & 800 & 1,600 & 3,200 \\
{}brute\_\allowbreak{}force & 450 & 800 & 1,600 & 3,200 \\
{}infiltration & 450 & 800 & 1,600 & 3,200 \\
{}web\_\allowbreak{}attacks & 450 & 450 & 450 & 450 \\
{}전체 학습/\allowbreak{}context & 2,700 & 4,450 & 8,450 & 16,450 \\
\bottomrule
\end{longtable}
\endgroup
\reportfigure{figures/exp87_extension_cic2018.pdf}{CIC2018: seed43의 기존 균형 끝점과 추가 확대 조건; 각 점에 F1 표시}


ToN-IoT의 실제 클래스별 제공 수다. 첫 열은 기존 균형 조건의 seed43 연결점(n=\allowbreak{}2,151)이며, 행은 고정 공격 test 표본 수 내림차순이다.

\begingroup
\footnotesize
\begin{longtable}{>{\raggedright\arraybackslash}p{\dimexpr 0.31618\linewidth-2\tabcolsep\relax}>{\raggedleft\arraybackslash}p{\dimexpr 0.22794\linewidth-2\tabcolsep\relax}>{\raggedleft\arraybackslash}p{\dimexpr 0.22794\linewidth-2\tabcolsep\relax}>{\raggedleft\arraybackslash}p{\dimexpr 0.22794\linewidth-2\tabcolsep\relax}}
\toprule
공격 클래스 & 기존 2151 & 3200 & 6400 \\ \midrule
\endfirsthead
공격 클래스 & 기존 2151 & 3200 & 6400 \\ \midrule
\endhead
{}xss & 2,151 & 3,200 & 6,400 \\
{}password & 2,151 & 3,200 & 6,400 \\
{}ddos & 2,151 & 3,200 & 6,400 \\
{}injection & 2,151 & 3,200 & 6,400 \\
{}backdoor & 2,151 & 3,200 & 6,400 \\
{}dos & 2,151 & 3,200 & 6,400 \\
{}scanning & 2,151 & 3,200 & 6,400 \\
{}mitm & 2,151 & 3,200 & 3,285 \\
{}ransomware & 2,151 & 2,151 & 2,151 \\
{}전체 학습/\allowbreak{}context & 19,359 & 27,751 & 50,236 \\
\bottomrule
\end{longtable}
\endgroup
\reportfigure{figures/exp87_extension_toniot.pdf}{ToN-IoT: seed43의 기존 균형 끝점과 추가 확대 조건; 각 점에 F1 표시}


\begingroup
\footnotesize
\begin{longtable}{>{\raggedright\arraybackslash}p{\dimexpr 0.14583\linewidth-2\tabcolsep\relax}>{\raggedleft\arraybackslash}p{\dimexpr 0.14583\linewidth-2\tabcolsep\relax}>{\raggedleft\arraybackslash}p{\dimexpr 0.14583\linewidth-2\tabcolsep\relax}>{\raggedleft\arraybackslash}p{\dimexpr 0.18155\linewidth-2\tabcolsep\relax}>{\raggedleft\arraybackslash}p{\dimexpr 0.18155\linewidth-2\tabcolsep\relax}>{\raggedleft\arraybackslash}p{\dimexpr 0.19940\linewidth-2\tabcolsep\relax}}
\toprule
데이터 & 목표 n & 총 사례 & XGB F1 (\%) & PFN F1 (\%) & PFN−XGB (\%p) \\ \midrule
\endfirsthead
데이터 & 목표 n & 총 사례 & XGB F1 (\%) & PFN F1 (\%) & PFN−XGB (\%p) \\ \midrule
\endhead
{}CIC2018 & 450 & 2,700 & 89.55 & 85.48 & -4.08 \\
{}CIC2018 & 800 & 4,450 & 90.43 & 86.31 & -4.12 \\
{}CIC2018 & 1,600 & 8,450 & 93.04 & 86.42 & -6.62 \\
{}CIC2018 & 3,200 & 16,450 & 92.98 & 86.67 & -6.31 \\
{}ToN-IoT & 2,151 & 19,359 & 92.27 & 92.78 & +\allowbreak{}0.50 \\
{}ToN-IoT & 3,200 & 27,751 & 93.06 & 93.03 & -0.03 \\
{}ToN-IoT & 6,400 & 50,236 & 94.47 & 92.00 & -2.47 \\
\bottomrule
\end{longtable}
\endgroup
CIC2018 XGB은 기존 균형 끝점 대비 현재 최대 목표 n=\allowbreak{}3,200에서 89.55 \ensuremath{\rightarrow} 92.98\% (+\allowbreak{}3.43\%p)다. 현재는 seed43 관찰이며 다른 seed의 재현 여부는 아직 확인하지 않았다.

CIC2018 TabPFN은 기존 균형 끝점 대비 현재 최대 목표 n=\allowbreak{}3,200에서 85.48 \ensuremath{\rightarrow} 86.67\% (+\allowbreak{}1.19\%p)다. 현재는 seed43 관찰이며 다른 seed의 재현 여부는 아직 확인하지 않았다.

ToN-IoT XGB은 기존 균형 끝점 대비 현재 최대 목표 n=\allowbreak{}6,400에서 92.27 \ensuremath{\rightarrow} 94.47\% (+\allowbreak{}2.20\%p)다. 현재는 seed43 관찰이며 다른 seed의 재현 여부는 아직 확인하지 않았다.

ToN-IoT TabPFN은 기존 균형 끝점 대비 현재 최대 목표 n=\allowbreak{}6,400에서 92.78 \ensuremath{\rightarrow} 92.00\% (-0.78\%p)다. 현재는 seed43 관찰이며 다른 seed의 재현 여부는 아직 확인하지 않았다.

CIC2018 web\_\allowbreak{}attacks와 ToN-IoT ransomware는 기존 끝점에서 이미 가용 학습 사례를 모두 사용했으므로, 이 두 클래스의 변화는 새로운 해당 공격 행을 추가한 효과가 아니다. ToN-IoT mitm은 2,151개에서 최대 3,285개로 실제 사례도 추가한다. 아래는 기존 균형 끝점과 가장 큰 완료 조건의 seed43 precision·recall·F1이며, 다른 공격 사례의 증가 및 클래스 비율 변화와 함께 해석한다.

\begingroup
\footnotesize
\begin{longtable}{>{\raggedright\arraybackslash}p{\dimexpr 0.24510\linewidth-2\tabcolsep\relax}>{\raggedright\arraybackslash}p{\dimexpr 0.13480\linewidth-2\tabcolsep\relax}>{\raggedleft\arraybackslash}p{\dimexpr 0.13480\linewidth-2\tabcolsep\relax}>{\raggedleft\arraybackslash}p{\dimexpr 0.18627\linewidth-2\tabcolsep\relax}>{\raggedleft\arraybackslash}p{\dimexpr 0.16422\linewidth-2\tabcolsep\relax}>{\raggedleft\arraybackslash}p{\dimexpr 0.13480\linewidth-2\tabcolsep\relax}}
\toprule
데이터 /\allowbreak{} 클래스 & 모델 & 목표 n & Precision (\%) & Recall (\%) & F1 (\%) \\ \midrule
\endfirsthead
데이터 /\allowbreak{} 클래스 & 모델 & 목표 n & Precision (\%) & Recall (\%) & F1 (\%) \\ \midrule
\endhead
{}CIC2018 /\allowbreak{} web\_\allowbreak{}attacks & XGB & 450 & 34.06 & 98.43 & 50.61 \\
{}CIC2018 /\allowbreak{} web\_\allowbreak{}attacks & XGB & 3,200 & 44.87 & 92.91 & 60.51 \\
{}CIC2018 /\allowbreak{} web\_\allowbreak{}attacks & TabPFN & 450 & 13.35 & 99.21 & 23.53 \\
{}CIC2018 /\allowbreak{} web\_\allowbreak{}attacks & TabPFN & 3,200 & 12.81 & 98.43 & 22.67 \\
{}ToN-IoT /\allowbreak{} mitm & XGB & 2,151 & 37.98 & 99.05 & 54.91 \\
{}ToN-IoT /\allowbreak{} mitm & XGB & 6,400 & 56.60 & 95.61 & 71.11 \\
{}ToN-IoT /\allowbreak{} mitm & TabPFN & 2,151 & 34.78 & 99.57 & 51.55 \\
{}ToN-IoT /\allowbreak{} mitm & TabPFN & 6,400 & 37.67 & 98.19 & 54.45 \\
{}ToN-IoT /\allowbreak{} ransomware & XGB & 2,151 & 97.52 & 99.86 & 98.68 \\
{}ToN-IoT /\allowbreak{} ransomware & XGB & 6,400 & 99.58 & 99.58 & 99.58 \\
{}ToN-IoT /\allowbreak{} ransomware & TabPFN & 2,151 & 93.65 & 99.86 & 96.66 \\
{}ToN-IoT /\allowbreak{} ransomware & TabPFN & 6,400 & 99.02 & 99.29 & 99.15 \\
{}ToN-IoT /\allowbreak{} dos & XGB & 2,151 & 85.55 & 99.16 & 91.86 \\
{}ToN-IoT /\allowbreak{} dos & XGB & 6,400 & 89.44 & 99.18 & 94.06 \\
{}ToN-IoT /\allowbreak{} dos & TabPFN & 2,151 & 96.63 & 98.70 & 97.65 \\
{}ToN-IoT /\allowbreak{} dos & TabPFN & 6,400 & 78.99 & 99.16 & 87.93 \\
\bottomrule
\end{longtable}
\endgroup
ToN-IoT TabPFN의 DoS precision은 96.63 \ensuremath{\rightarrow} 78.99\%, recall은 98.70 \ensuremath{\rightarrow} 99.16\%다. 같은 고정 test에서 XSS를 DoS로 잘못 분류한 행은 516 \ensuremath{\rightarrow} 5,490개다. 표본 확대 후 품질 변화를 공격 누락과 다른 공격의 오분류 유입으로 나누어 확인해야 한다.

다음 확인은 끝점·직전 조건의 다른 seed 재현성과, 크게 늘어난 오분류 쌍에 대한 context 구성 분석이다. 단일 seed 결과만으로 최적 context 크기나 모델의 최종 성능 상한을 정하지 않는다. 표본 수와 클래스 비율의 영향을 분리하려면 별도 배분 대조가 필요하며 이번 결과만으로 원인을 확정하지 않는다.

추가 실행의 입력·예측·비용과 검증 근거는 \texttt{tabpfn/\allowbreak{}\allowbreak{}r\allowbreak{}esults/\allowbreak{}\allowbreak{}v\allowbreak{}7/\allowbreak{}\allowbreak{}exp87/\allowbreak{}\allowbreak{}20261022\allowbreak{}\_\allowbreak{}\allowbreak{}164549\_\allowbreak{}\allowbreak{}class\_\allowbreak{}\allowbreak{}ca\allowbreak{}pped\_\allowbreak{}\allowbreak{}ext\allowbreak{}ension\_\allowbreak{}\allowbreak{}s\allowbreak{}eed43\_\allowbreak{}\allowbreak{}43\allowbreak{}f6dd6c}에 보존한다. 비용은 values.csv의 preparation\_\allowbreak{}seconds와 inference\_\allowbreak{}seconds로 분리했다.

\subsection{재현·검증 자료}
실행 경로 \texttt{tabpfn/\allowbreak{}\allowbreak{}r\allowbreak{}esults/\allowbreak{}\allowbreak{}v\allowbreak{}7/\allowbreak{}\allowbreak{}exp87/\allowbreak{}\allowbreak{}20261022\allowbreak{}\_\allowbreak{}\allowbreak{}163739\_\allowbreak{}\allowbreak{}full\_\allowbreak{}\allowbreak{}lea\allowbreak{}rning\_\allowbreak{}\allowbreak{}cu\allowbreak{}rve\_\allowbreak{}\allowbreak{}five\allowbreak{}\_\allowbreak{}\allowbreak{}seeds\_\allowbreak{}\allowbreak{}c\allowbreak{}59213f5}의 config/\allowbreak{}protocol, source\_\allowbreak{}dependencies, data, jobs, events/\allowbreak{}decisions, CSV를 확인한다. 예측 혼동행렬과 모델 간 입력 ID를 재검증했다. 원고·슬라이드는 자동 변경하지 않는다.

그림에 표시한 평균·표준편차와 완료 n은 실행 폴더의 figure\_\allowbreak{}values.json에 저장했다. 비용과 클래스별 결과도 동일 seed 집합만 사용한다.


\end{document}
